\section{Persistent-State Serialization and Compute Scope}
\label{sec:embedded}

This appendix makes every byte in the budget claim reproducible and separates
what the budget covers from what it does not.

\subsection{What the 16\,KiB covers}

\textbf{The 16\,KiB constraint applies to persistent adaptation state, not to
peak training SRAM.} Table~\ref{tab:memcats} lists every memory category and its
status. The included categories are all quantities a device must retain between
adaptation steps; the excluded categories are transient or implementation
specific, and none of them is measured in this work.

\begin{table}[H]
\centering
\small
\caption{Memory categories and their treatment. ``Replication only'' marks the
implementation reserve, which is held back in the reserve replication of
Section~\ref{sec:reserve} but \emph{not} in the primary configuration. No
sentence in this paper claims that total training SRAM is 16\,KiB.}
\label{tab:memcats}
\begin{tabular}{lccl}
\toprule
Memory category & In the 16\,KiB? & Persistent? & Status \\
\midrule
Trainable weights          & yes & yes & analytical \\
Gradients                  & yes & assumption-dependent & analytical \\
Optimizer moments          & yes & yes & analytical \\
Replay values              & yes & yes & analytical \\
Replay labels              & yes & yes & analytical \\
Buffer and quantization metadata & yes & yes & analytical \\
Alignment padding          & yes & yes & analytical \\
Implementation reserve     & replication only & yes & analytical \\
\midrule
Forward activations        & no  & no  & unmeasured \\
Backward activations       & no  & no  & unmeasured \\
Stack and workspace        & no  & no  & unmeasured \\
DMA and allocator overhead & no  & no  & unmeasured \\
Flash write endurance      & no  & --- & unmeasured \\
\bottomrule
\end{tabular}
\end{table}

Gradients are marked assumption-dependent because an implementation that applies
updates immediately per batch need not retain a full gradient vector across
steps. We include them so the accounting is conservative.

\subsection{Packed replay records}

Replay entries are serialized as fixed-stride packed records. With INT8 replay
values, the raw record is

\begin{verbatim}
typedef struct __attribute__((packed)) {
    int8_t  ecg[198];   /* R-peak-centred beat, INT8            */
    int8_t  pre_rr;     /* normalised previous RR interval      */
    int8_t  post_rr;    /* normalised following RR interval     */
    uint8_t label;      /* AAMI class 0..4                      */
    uint8_t valid;      /* slot occupancy flag                  */
    uint8_t class_id;   /* selector bookkeeping                 */
} replay_raw_t;         /* 203 B packed                          */
\end{verbatim}

and the post-pool record is

\begin{verbatim}
typedef struct __attribute__((packed)) {
    int8_t  pooled[16]; /* post-pool embedding, INT8            */
    int8_t  pre_rr;
    int8_t  post_rr;
    uint8_t label;
    uint8_t valid;
    uint8_t class_id;
} replay_postpool_t;    /* 21 B packed                           */
\end{verbatim}

Both structures are all-\texttt{int8}/\texttt{uint8}, so the packed size, the
natural ABI size, and the 1-byte-aligned stride coincide: 203\,B and 21\,B
respectively. Under 4-byte alignment the strides become 204\,B and 24\,B, and
under 8-byte alignment 208\,B and 24\,B; we report the 1-byte figures and
account for the padding explicitly rather than assuming a favorable ABI.

Beyond the per-record fields, each buffer carries fixed metadata: a 4-byte write
index, a 4-byte occupancy count, 16 bytes of reservoir/RNG state, one 4-byte
counter per AAMI class (20\,B), and per-quantized-tensor scale and zero point
(4\,B and 1\,B each). This totals 49\,B for the post-pool buffer, which quantizes
one tensor, and 54\,B for the raw buffer, which quantizes two (the beat and the
RR pair). Table~\ref{tab:arm_definitions} reports the resulting totals.

\subsection{Implementation reserve}

A configuration that consumes 16{,}383 of 16{,}384 bytes is not deployable: it
leaves nothing for a firmware version field, a checksum, or a future field
addition. A deployable rule would hold back a reserve,
\[
B_{\mathrm{used}} \;\le\; B_{\mathrm{budget}} - B_{\mathrm{reserve}},
\]
which at $B_{\mathrm{reserve}} = 1024$\,B lowers the effective ceiling to
15{,}360\,B.

We are explicit about how this relates to the results, because the two are easy
to conflate. \textbf{The primary analysis does \emph{not} apply a reserve.} Its
configurations are generated against the full 16{,}384\,B ceiling, because that
is what the pre-specified analysis plan fixed, and changing the byte rule after
the fact would forfeit the pre-specification. The reserve is instead reported as
a \emph{separate replication} (Section~\ref{sec:reserve},
Table~\ref{tab:reserve}): the entire A0--A5 grid re-run at the 15{,}360\,B
ceiling with regenerated replay counts, which preserves every significant gain
over the frozen model but \emph{not} the exact arm ordering --- A1 rises to
$0.793$ and A5 falls to $0.792$, so A4 leads under the reserve where A5 leads
without it. Table~\ref{tab:arm_definitions} lists both count sets side
by side so no arm's performance is ever quoted against a byte figure it was not
measured under. A deployment should use the reserved counts; the primary table
documents the pre-specified configuration.

\subsection{Compute is not matched across arms}

The arms are matched on persistent bytes but \emph{not} on compute, and the
direction of the gap is unambiguous even though this paper does not quantify it.
Raw replay requires re-forwarding stored beats through the encoder and
propagating gradients through the LoRA training path, whereas post-pool replay
enters the graph after the encoder and bypasses it entirely. Encoder adaptation
therefore carries a substantially greater compute and energy cost per adaptation
step than head-only or post-pool replay.

\textbf{We do not report a numeric ratio.} An earlier version of this appendix
gave analytical MAC totals and a fixed multiplier, charging backward cost as a
multiple of the forward cost of the trainable path only. That convention is
defensible for the head and post-pool arms, whose trainable parameters sit after
the frozen encoder, but it is not a safe basis for a headline number in the
encoder-LoRA arms: reverse-mode differentiation must propagate activation
gradients \emph{through} frozen operations to reach LoRA factors inserted inside
the encoder, so a backward pass cannot be priced from the trainable parameter
count alone. Rather than publish a ratio whose derivation we cannot fully stand
behind, we state the qualitative asymmetry and leave the magnitude to a
measured study --- a layer-by-layer forward and backward derivation, or a
profiler trace on the target runtime --- in the journal version.

This matters for a deployment decision and not for the allocation result. The
memory-allocation contribution of this paper is an accounting of persistent
\emph{bytes}, and every comparison in it is byte-matched by construction; none
of the conclusions depends on the compute figures withdrawn here. A reader
choosing between these arms on a duty-cycled device should nevertheless treat
encoder adaptation as the more expensive option and measure it before
committing.
